\documentclass{article}
\usepackage{fontspec}
\setmainfont{Times New Roman}
\usepackage{amsmath}
\usepackage{amssymb}
\usepackage{amsfonts}
\usepackage{geometry}
\geometry{a4paper, margin=1in}

\title{Лекция 1: Въведение в теория на вероятностите. Събития и сигма-алгебри}
\author{}
\date{}

\begin{document}

\maketitle

\section{Въведение и организационни въпроси}
Курсът се състои от три основни компонента: лекции, семинарни занятия и практикум. Оценката се формира като средно претеглена от тези компоненти, като за успешно преминаване е необходим минимум от 3.00 на всеки от тях. Оценките от семинарните занятия и практикума важат до септемврийската поправителна сесия. Целта на курса е изграждане на стохастично мислене, което е фундаментално за съвременни научни дисциплини като машинното обучение (machine learning) и изкуствения интелект.

\section{Мотивация. Ролята на стохастиката}
Стохастичният апарат предоставя мощен инструмент за моделиране на явления, които не са напълно детерминирани. Няколко примера от историята и ежедневието илюстрират това:
\begin{itemize}
    \item \textbf{Брауново движение (Робърт Браун 1827 г., Алберт Айнщайн 1905 г.):} Хаотичното движение на поленова частица във вода е резултат от ударите на водните молекули. Траекторията ѝ има безкрайна дължина и е никъде диференцируема, но служи като отлично математическо приближение (Жан Перен използва това за оценка на числото на Авогадро). Днес подобни стохастични процеси се използват широко във финансовата математика (например моделът на Блек и Шоулс).
    \item \textbf{Тото 6 от 49:} На 6 септември 2009 г. в България се падат едни и същи числа в два последователни тиража на тотото. Използвайки теория на вероятностите, може да се покаже, че при общо 10 000 изиграни тиража подобно събитие не е толкова невероятно, колкото изглежда интуитивно, и не е непременно индикация за измама.
\end{itemize}

\section{Случайни експерименти и събития}
\textbf{Дефиниция:} Под \emph{случаен експеримент} разбираме изследване или опит, чийто изход не е предварително предопределен.

\textbf{Дефиниция:} За даден случаен експеримент, множеството от всички възможни елементарни изходи се нарича множество от елементарни събития и се обозначава с $\Omega$. Отделните елементарни събития се обозначават с $\omega$.

\textbf{Примери:}
\begin{enumerate}
    \item Хвърляне на монета: Изходът е бинарен.
    \[ \Omega = \{ \text{ези}, \text{тура} \} \quad \text{или} \quad \Omega = \{ 0, 1 \} \]
    \item Хвърляне на зарче: Възможните изходи са шест.
    \[ \Omega = \{ 1, 2, 3, 4, 5, 6 \} \]
    Пример за елементарно събитие е падането на петица: $\omega_5 = \{ 5 \}$.
    \item Тото 6 от 49: Множеството от всички възможни изтеглени шестици в един тираж има размерност:
    \[ N = \binom{49}{6} = 13\,983\,816 \]
    Тогава $\Omega = \{ \omega_1, \omega_2, \dots, \omega_N \}$.
    \item Време на живот на електронен компонент (напр. процесор):
    \[ \Omega = \mathbb{R}_+ = [0, \infty) \]
\end{enumerate}

\section{Операции със събития}
Събитията се разглеждат като подмножества на $\Omega$.

\textbf{Дефиниция (Допълнение):} Нека $A \subseteq \Omega$. Допълнението на $A$ се дефинира като:
\[ A^c = \Omega \setminus A \]
Това означава, че ако $\omega \in A^c$, то $\omega \notin A$. Също така е в сила:
\[ \Omega = A \cup A^c \quad \text{и} \quad (A^c)^c = A \]
\emph{Примери:} При хвърляне на зарче, ако $A = \{ 1, 3, 5 \}$ (нечетно число), то $A^c = \{ 2, 4, 6 \}$ (четно число). Друг пример: ако $A = \{ \text{жени} \}$, то $A^c = \{ \text{мъже} \}$.

\subsection*{Свойства на операциите}
За произволни събития $A, B$ и $C$ са в сила следните свойства:
\begin{itemize}
    \item \textbf{a) Комутативност:}
    \[ A \cup B = B \cup A; \quad A \cap B = B \cap A \]
    \item \textbf{б) Асоциативност:}
    \[ A \cup (B \cup C) = (A \cup B) \cup C = A \cup B \cup C \]
    \[ A \cap (B \cap C) = (A \cap B) \cap C = A \cap B \cap C \]
    \item \textbf{в) Дистрибутивност:}
    \[ A \cup (B \cap C) = (A \cup B) \cap (A \cup C) \]
    \item \textbf{г) Безкрайни операции и закони на Де Морган:}
    Ако $A_i, i \ge 1$ са събития, то обединението им се дефинира като множеството от елементарни събития, принадлежащи на поне едно от тях:
    \[ \bigcup_{i=1}^{\infty} A_i = \{ \omega \in \Omega : \omega \in A_i \text{ за поне едно } i \} \]
    Сечението им се дефинира като множеството от елементарни събития, принадлежащи на всяко едно от тях:
    \[ \bigcap_{i=1}^{\infty} A_i = \{ \omega \in \Omega : \omega \in A_i \text{ за всяко } i \} \]
    Връзката между тези операции се дава от законите на Де Морган:
    \[ (A \cup B)^c = A^c \cap B^c \]
    \[ \left( \bigcup_{i=1}^{\infty} A_i \right)^c = \bigcap_{i=1}^{\infty} A_i^c \]
    \[ \left( \bigcap_{i=1}^{\infty} A_i \right)^c = \bigcup_{i=1}^{\infty} A_i^c \]
\end{itemize}

\section{Сигма-алгебра}
За да можем да дефинираме коректно вероятност, ни е необходима структура върху колекцията от събития, които ще измерваме.

\textbf{Дефиниция ($\sigma$-алгебра):} Нека $\Omega$ е множество от елементарни събития, а $\mathcal{A}$ е колекция от подмножества на $\Omega$. Тогава $\mathcal{A}$ се нарича $\sigma$-алгебра, ако са изпълнени следните условия:
\begin{enumerate}
    \item[i)] $\emptyset \in \mathcal{A}$
    \item[ii)] Ако $A \in \mathcal{A}$, то $A^c \in \mathcal{A}$
    \item[iii)] Ако $A, B \in \mathcal{A}$, то $A \cup B \in \mathcal{A}$. В по-общия случай, ако $A_i \in \mathcal{A}$ за $\forall i \ge 1$, то $\bigcup_{i=1}^{\infty} A_i \in \mathcal{A}$.
\end{enumerate}

\textbf{Следствие:} Нека $\mathcal{A}$ е $\sigma$-алгебра и $A_i \in \mathcal{A}$ за $\forall i \ge 1$. Тогава:
\[ B = \bigcap_{i=1}^{\infty} A_i \in \mathcal{A} \]
\emph{Доказателство:}
По законите на Де Морган имаме:
\[ B^c = \left( \bigcap_{i=1}^{\infty} A_i \right)^c = \bigcup_{i=1}^{\infty} A_i^c \]
От условие ii) следва, че $A_i^c \in \mathcal{A}$ за всяко $i$. Тогава от условие iii) безкрайното обединение на тези допълнения също е в $\mathcal{A}$, т.е. $B^c \in \mathcal{A}$. Прилагайки отново условие ii), получаваме:
\[ B = (B^c)^c \in \mathcal{A} \]

\subsection*{Примери за $\sigma$-алгебри}
\begin{enumerate}
    \item \textbf{Тривиална $\sigma$-алгебра:} $\mathcal{A} = \{ \emptyset, \Omega \}$.
    \item \textbf{Булеан (всички подмножества):} При хвърляне на монета $\Omega = \{ 0, 1 \}$, $\mathcal{A} = 2^\Omega$. 
    В експеримент с 3 възможни изхода (напр. попадане на стрела в една от три цветни области), $\Omega = \{ \omega_1, \omega_2, \omega_3 \}$. Броят на елементите в $\sigma$-алгебрата е $2^3 = 8$:
    \[ \mathcal{A} = 2^\Omega = \{ \emptyset, \Omega, \{ \omega_1 \}, \{ \omega_2 \}, \{ \omega_3 \}, \{ \omega_1, \omega_2 \}, \{ \omega_1, \omega_3 \}, \{ \omega_2, \omega_3 \} \} \]
    \item \textbf{Борелова $\sigma$-алгебра:} При неизброими пространства като $\Omega = \mathbb{R}_+ = [0, \infty)$, множеството от всички подмножества $2^\Omega$ е твърде голямо за коректно дефиниране на вероятност. Вместо това се използва Бореловата $\sigma$-алгебра, генерирана от интервали от вида $A = (a, b)$.
\end{enumerate}

\section{Формализиране на примера с Тото 6 от 49}
Нека разгледаме $10\,000$ тиража на тотото. Пространството от всички възможни изходи в тези тиражи е:
\[ \Omega = \bigotimes_{i=1}^{10\,000} \Omega_i \]
където $\Omega_i$ е множеството от всички възможни $N = 13\,983\,816$ шестици. Мощността на $\Omega$ е огромна:
\[ |\Omega| = (13\,983\,816)^{10\,000} \]
Дефинираме събитието:
\[ A = \{ \text{паднали са се едни и същи числа в два последователни тиража от } 10\,000 \} \]

За да запишем това математически, нека $A_i$ бъде събитието, при което изтеглената шестица в $i$-тия тираж съвпада с тази в $(i+1)$-вия тираж:
\[ A_i = \{ \omega^{(i)} = \omega^{(i+1)} \} \subseteq \Omega \]
Тогава търсеното събитие $A$ може да се представи като обединение на събитията $A_i$ за всеки два последователни тиража:
\[ A = \bigcup_{i=1}^{9999} A_i \]
Чрез тези дефиниции словесното описание се свежда до строг теоретико-вероятностен модел, който подлежи на количествен анализ.

\section{Парадоксът с двата плика}
Като задача за размисъл се разглежда следният парадокс:
Имате два плика, съдържащи числа $a$ и $b$, като $a < b$. Избирате един от тях напълно случайно, отваряте го и виждате числото. Има ли стратегия, която да ви гарантира вероятност строго по-голяма от $1/2$ да изберете плика с по-голямата сума?
\end{document}